\begin{abstract}
What can a synthetic fraud benchmark show, and what can it not? We answer the question for
FraudShield-EAC, a seeded synthetic benchmark of East African digital payments (mobile money, USSD,
agent banking, card, online and bank transfer, across \nCountries{} countries). The benchmark is easy
in a specific way: a single velocity feature separates fraud from legitimate payments at an AUC of
\nGateFloor{} on its own, so we report every ranking metric as a margin over that feature rather than
over chance.

Trained on all \nFvTrainRows{} training rows, a calibrated XGBoost and LightGBM ensemble reaches an AUC
of \nFvAUC{}, a margin of \nFvMargin{}. Its recall at the specified flag threshold, \nFvRecallSix{},
misses the target of \nTargetRecall{}, and more training data did not close that gap in any run we
measured. At a fixed \nTrainUsed{} rows, the corpus they were drawn from moved mean recall across a
range of \nSelMeanRange{}, against \nFvSeedGainFixed{} from \nTrainUsed{} rows to the whole period on one
corpus; no interval was computed for these differences. Ablations measure
redundancy rather than importance, and leave-one-country-out and a fraud variant held out in time are
both null. Refitting without one fraud type at a time costs little recall at a \nFPRBudget{} false-positive
budget for velocity, SIM-swap and
mule fraud and much more for agent fraud and synthetic identity.

The one failure of the production model on a pre-specified variant is a variant with no burst
structure, modelled on a documented Kenyan mobile-money scam and committed, with the diagnostic model's predicted recall,
alongside the generator code that produces it, before its evidence. It is never caught at the flag threshold; at
a \nFPRBudget{} false-positive budget the model trained on \nTrainUsed{} rows catches \nGateRevNewBudgetCaught{} of the
\nGateRevNewBudgetRows{} rows whose recipient the account had never paid and none of the
\nGateRevKnownBudgetRows{} sent to a known payee, small counts from an observational split that do not isolate a cause. All results
are on a synthetic benchmark whose fraud parameters are modelling choices; we state what a validation
on real data would need to show.
\end{abstract}
